%% kapa: exam=1 round=22 year=2011 period=1 form=A subject=civil_law subject_ko=민법 seq=11 number=11
%% source: Q-Net 감정평가사 기출문제 (cst00309) artlSeq=5036790 file=extracted/1008/5036790/2011년 제22회 감정평가사 1차시험 기출문제/1교시.hwp
%% conversion: kordoc 4.24.51 -> tools/gen1.py (auto); review=proofread-sonnet; flags=text-verified-hwp
\begin{question}{11}
관습법에 관한 설명으로 옳은 것은? (다툼이 있으면 판례에 의함)
\begin{choices}
  \item 온천수에 관한 권리는 관습법상 물권이 아니라 상린관계상 공용수 또는 생활상 필요한 용수에 관한 권리에 해당한다.
  \item 미등기 무허가건물의 양수인에게 소유권에 준하는 관습법상의 물권이 인정된다.
  \item 미등기건물을 대지와 함께 매도하였으나 대지에 관하여만 매수인 앞으로 소유권이전등기가 경료된 경우, 매도인에게 관습상의 법정지상권이 인정된다.
  \item 일반 주민들이 다른 사람의 공동 사용을 방해하지 않는 한 자유로이 이용할 수 있는 공원이용권은 관습법상 물권이다.
  \item 도로가 개설된 후 장기간에 걸쳐 일반의 통행에 제공되어 왔고 공로로 통하는 유일한 것일지라도 관습법상 통행권이 인정되지 않는다.
\end{choices}
\end{question}
\answer{5}
